\documentclass[11pt]{article}
\usepackage[margin=1in]{geometry}
\usepackage{amsmath,amssymb,amsthm,mathtools}
\usepackage[hidelinks]{hyperref}

\newtheorem{theorem}{Theorem}[section]
\newtheorem{lemma}[theorem]{Lemma}
\newtheorem{proposition}[theorem]{Proposition}
\theoremstyle{remark}
\newtheorem{remark}[theorem]{Remark}
\newcommand{\splus}{s^+}
\newcommand{\sminus}{s^-}
\newcommand{\Arg}{\operatorname{Arg}}

\title{Strict Positive Square Energy of Every Bicyclic Graph}
\author{AI-generated manuscript; no human author is claimed}
\date{30 July 2026}

\begin{document}
\maketitle

\begin{abstract}
Let $G$ be a finite simple connected graph with $n$ vertices and $n+1$
edges.  We prove the strict inequality $\splus(G)>n$, where $\splus$ is the
sum of the squares of the positive adjacency eigenvalues.  The $2$-core is
either a bicyclic cactus (a figure-eight or handcuff/dumbbell) or a theta;
arbitrary rooted trees may be attached at its vertices.  The cactus class is
covered by the strict theorem for every nonunicyclic cactus.  For the
non-cactus class we give the analytic attached-theta argument: tree messages,
an open-right-half-plane Sachs expansion, and exact endpoint packets controlling
zero, one, or two hostile odd cycles.  The delicate negative Sachs coefficient
of a $0\pmod4$ cycle is handled by explicit factorizations.  The phase charge
theorem gives, for every nonbipartite attached theta,
\[
 \splus(G)-n\geq5-2\sqrt5>0.
\]
The cited finite verifier checks packet identities and hostile mutations, but
the proof is uniform and does not use a finite path-length census.
\end{abstract}

\section{Statement and structural reduction}

All graphs are finite, simple, and undirected.  For the adjacency eigenvalues
$\lambda_1,\ldots,\lambda_n$, set
\[
 \splus(G)=\sum_{\lambda_j>0}\lambda_j^2,\qquad
 \sminus(G)=\sum_{\lambda_j<0}\lambda_j^2,\qquad
 D(G)=\splus(G)-\sminus(G).
\]
Since $\operatorname{tr}A(G)^2=2|E(G)|$, a bicyclic graph satisfies
\begin{equation}\label{eq:energy}
 \splus(G)+\sminus(G)=2n+2,
 \qquad \splus(G)=n+1+\frac12D(G).
\end{equation}

\begin{theorem}\label{thm:main}
Every finite simple connected graph $G$ with $|V(G)|=n$ and
$|E(G)|=n+1$ satisfies
\[
 \boxed{\splus(G)>n}.
\]
\end{theorem}

Recall that the $2$-core is obtained by repeatedly deleting vertices of degree
at most one.

\begin{proposition}[Bicyclic core dichotomy]\label{prop:structure}
The $2$-core $H$ of a finite simple connected bicyclic graph is exactly one of
the following:
\begin{enumerate}
 \item a cactus with two cyclic blocks, either sharing one cut vertex
 (figure-eight or tight handcuff), or vertex-disjoint and joined by a path
 (dumbbell or loose handcuff);
 \item a theta graph, consisting of three internally vertex-disjoint paths of
 positive lengths between two distinct terminals, with at most one path of
 length one.
\end{enumerate}
Every component outside $H$ is a tree attached to exactly one vertex of $H$.
\end{proposition}

\begin{proof}
The core is connected, has minimum degree at least two, and retains
$|E(H)|-|V(H)|=1$.  Thus
\[
 \sum_{v\in V(H)}(\deg_H(v)-2)=2.
\]
Either one vertex has degree four and all others degree two, or two vertices
have degree three and all others degree two.  The first case gives two cycles
meeting at the exceptional vertex.  In the second case suppress maximal
degree-two paths.  If a suppressed edge is a bridge, the remaining half-edges
form one cycle at each endpoint and the bridge expands to the joining path.
Otherwise the suppressed multigraph has three parallel terminal-to-terminal
edges, which expand to a theta.  Simplicity forbids two direct terminal edges.
Reversing leaf deletion attaches rooted trees, each at one core vertex.
\end{proof}

The cactus alternative follows from the following strict theorem, whose
manuscript includes the rank-uniform DNN frontier and closure of its residual
families.

\begin{theorem}[Nonunicyclic cactus theorem \cite{Cactus}]\label{thm:cactus}
Every finite simple connected cactus of cyclomatic rank at least two satisfies
$\splus(G)>|V(G)|$.
\end{theorem}

We prove the attached-theta alternative in enough detail that the main result
does not treat the research note as a black box.

\section{Tree messages and the right half-plane}

For positive activities $a=(a_v)$ and a graph $J$, define
\begin{equation}\label{eq:Z}
 Z_J(a)=\sum_{M\text{ matching of }J}
             \prod_{v\text{ unmatched by }M}a_v.
\end{equation}
Orient each off-core branch toward its core root.  The leaf recurrence gives
\begin{equation}\label{eq:messages}
 Q_{u\to v}=t+\sum_{w\text{ child of }u}Q_{w\to u}^{-1}>0,
 \qquad
 a_v=t+\sum_{u\text{ off-core at }v}Q_{u\to v}^{-1}\geq t.
\end{equation}
Extracting branch partitions gives one factor $K(t)>0$.  It remains common
when a Sachs cycle deletes a core root: the detached branches contribute their
full partitions.  Thus no cycle-deletion term carries a hidden quotient.

Normalize
\begin{equation}\label{eq:Psi}
 \Psi_G(t)=i^{-n}\det(itI-A(G))=\prod_{j=1}^n(t+i\lambda_j),\qquad t>0.
\end{equation}
Its continuous argument tending to zero at infinity is
$\Theta_G(t)=\sum_j\arctan(\lambda_j/t)$, and the signed Coulson identity is
\begin{equation}\label{eq:coulson}
 D(G)=-\frac4\pi\int_0^\infty t\Theta_G(t)\,dt.
\end{equation}
This follows by integrating the identity for $\lambda|\lambda|$, using
$\sum_j\lambda_j=0$, and integrating by parts.

Let theta paths $P_0,P_1,P_2$ join terminals $x,y$, with
$\ell_j=|E(P_j)|$.  If their parities agree, the graph and all its tree
attachments are bipartite, so spectral symmetry gives
$\splus(G)=|E(G)|=n+1$.  Otherwise label the two odd and one even cycles
\begin{equation}\label{eq:cycles}
 C_1=P_0\cup P_1,\qquad C_2=P_0\cup P_2,
 \qquad C_e=P_1\cup P_2.
\end{equation}
Put
\[
 d_j=Z_{P_j-\{x,y\}}(a),\qquad
 \varepsilon=\begin{cases}+1,&|C_e|\equiv2\pmod4,\\
                           -1,&|C_e|\equiv0\pmod4.
             \end{cases}
\]
Call an odd cycle hostile when its length is $1\pmod4$, favorable when it is
$3\pmod4$, and set $\eta_j=+1$ or $-1$ accordingly.  Sachs expansion gives
\begin{equation}\label{eq:sachs-theta}
 \frac{\Psi_G(t)}{K(t)}=W=R+2i(\eta_1d_2+\eta_2d_1),
 \qquad R=Z_H(a)+2\varepsilon d_0.
\end{equation}
No two theta cycles are vertex-disjoint, so no double-cycle term occurs.

\begin{lemma}[Open right half-plane]\label{lem:rhp}
For every $t>0$, $R>0$.  Thus the continuous argument of $W$ is its principal
argument in $(-\pi/2,\pi/2)$.
\end{lemma}

\begin{proof}
Partition matchings first on the even cycle.  The omitted-path packet below
gives $Z_H(a)=z_ed_0+E_e$, where
$z_e=Z_{C_e}(a|_{C_e})$ and $E_e\geq0$.  The even cycle has two perfect
matchings and its empty matching has positive weight, so $z_e>2$.  Hence
\[
 R=(z_e+2)d_0+E_e>0\quad(\varepsilon=+1),\qquad
 R=(z_e-2)d_0+E_e>0\quad(\varepsilon=-1).
\]
\end{proof}

\section{The exact omitted-path packet}

For a list of positive numbers define the continuant
\[
 K()=1,\quad K(c_1)=c_1,\quad
 K(c_1,\ldots,c_r)=c_rK(c_1,\ldots,c_{r-1})+K(c_1,\ldots,c_{r-2}).
\]
It obeys
\begin{equation}\label{eq:continuant-det}
 K(c_1,\ldots,c_r)K(c_2,\ldots,c_{r-1})
 -K(c_1,\ldots,c_{r-1})K(c_2,\ldots,c_r)=(-1)^r,
\end{equation}
the determinant identity behind terminal path elimination.

Let $C$ be one theta cycle and let its omitted path be
$Q:x=v_0,v_1,\ldots,v_r=y$, with $c_j=a_{v_j}$.  Define
\begin{align*}
 D_Q&=K(c_1,\ldots,c_{r-1}),\\
 L_Q&=\begin{cases}0,&r=1,\\K(c_2,\ldots,c_{r-1}),&r\geq2,\end{cases}
 &R_Q&=\begin{cases}0,&r=1,\\K(c_1,\ldots,c_{r-2}),&r\geq2,\end{cases}\\
 M_Q&=\begin{cases}1,&r=1,\\0,&r=2,\\
 K(c_2,\ldots,c_{r-2}),&r\geq3.
 \end{cases}
\end{align*}
Empty continuants are one.  If $z_C=Z_C(a|_C)$ and superscripts indicate
deleted terminals, partitioning by the first and last edges of $Q$ gives the
exact four-state identity
\begin{equation}\label{eq:packet}
 Z_H(a)=z_CD_Q+z_C^xL_Q+z_C^yR_Q+z_C^{xy}M_Q.
\end{equation}
This includes a direct omitted edge ($r=1$) and a two-edge path ($r=2$), the
two cases in which informal empty-path conventions are most dangerous.

For a hostile cycle of length $q$, activity monotonicity gives
\begin{equation}\label{eq:weighted-cycle}
 z_C\geq Z_{C_q}(t)=:Z_q(t),
 \,\qquad
 0<\Arg(z_C+2i)\leq\Arg(Z_q(t)+2i).
\end{equation}
The comparator is the isolated weighted cycle after tree elimination, not the
bare attached theta; the latter comparison need not hold.

\section{Channel lemmas}

If neither odd cycle is hostile, \eqref{eq:sachs-theta} gives
$\operatorname{Im}W=-2(d_1+d_2)<0$, and hence $\Theta_G(t)<0$.  The one- and
two-hostile channels require exact control of the negative coefficient created
when $|C_e|\equiv0\pmod4$.

\begin{lemma}[One-hostile channel]\label{lem:one}
Suppose $C_1$ is hostile, $C_2$ is favorable, and $q=|C_1|$.  Then
\begin{equation}\label{eq:one-phase}
 \Theta_G(t)<\Arg(Z_q(t)+2i)\qquad(t>0).
\end{equation}
\end{lemma}

\begin{proof}
Set $B=d_2$, $A=d_1$, and $Z_q=Z_q(t)$, so
$\operatorname{Im}W=2(B-A)$.  If $|C_e|\equiv2\pmod4$, applying
\eqref{eq:packet} to $C_1$ gives
\begin{align}\label{eq:positive-one}
 R-Z_q(B-A)={}&(z_{C_1}-Z_q)B+Z_qA+z_{C_1}^xL_{P_2}\notag\\
 &+z_{C_1}^yR_{P_2}+z_{C_1}^{xy}M_{P_2}+2d_0>0.
\end{align}
If $B-A\leq0$, the phase is nonpositive.  Otherwise division by $RZ_q>0$
and Lemma~\ref{lem:rhp} give
$\arctan(2(B-A)/R)<\arctan(2/Z_q)$.

Suppose instead that $|C_e|\equiv0\pmod4$.  The last term in
\eqref{eq:positive-one} becomes $-2d_0$, so termwise positivity would be
invalid.  Define
\[
 F(a)=Z_H(a)-2d_0-Z_q(B-A).
\]
We show $F(a)\geq F(t)>0$ on $a_v\geq t$.  Let
$I_j=V(P_j)-\{x,y\}$.  Differentiating a matching partition deletes the
corresponding vertex, and hence
\[
 \partial_vF=\begin{cases}
 Z_{H-v},&v=x,y,\\
 Z_{H-v}+Z_qZ_{I_1-v},&v\in I_1,\\
 Z_{H-v}-Z_qZ_{I_2-v},&v\in I_2,\\
 Z_{H-v}-2Z_{I_0-v},&v\in I_0.
 \end{cases}
\]
For $v\in I_2$, unite a matching of $C_1$ and a matching of $I_2-v$, using
no boundary edge.  This is a coefficient-preserving injection into matchings
of $H-v$, so
$Z_{H-v}\geq z_{C_1}Z_{I_2-v}\geq Z_qZ_{I_2-v}$.  Similarly, using $C_e$ and
$I_0-v$ gives $Z_{H-v}>2Z_{I_0-v}$ because $Z_{C_e}(a)>2$.  Every derivative
is nonnegative, proving $F(a)\geq F(t)$.

For strict bare positivity put $a=\ell_0$, $b=\ell_1$, $c=\ell_2$, and let
$K_r=K(t,\ldots,t)$ ($r$ entries), with $K_{-1}=0$, $K_0=1$.  The residues
force $a$ even and $b,c$ odd.  A two-image, uncovered-vertex-preserving
matching injection gives coefficientwise
\begin{equation}\label{eq:bare-injection}
 (K_{a+b}+K_{a+b-2})K_{b-1}\geq2K_{a-1}.
\end{equation}
Indeed, for each matching of the interior of $P_0$, complete through $P_1$ in
the two parity choices and shift each retained $P_0$ edge by one in the second
image.  The images are disjoint because exactly the second uses the first
$P_0$ edge, and both add $b$ matched edges, preserving every power of $t$.
Formula \eqref{eq:packet} now yields
\[
 F(t)=Z_qK_{b-1}-2K_{a-1}+K_{a-1}K_{b-1}>0\quad(c=1),
\]
while for odd $c\geq3$ it yields
\[
 F(t)=Z_qK_{b-1}-2K_{a-1}
 +2K_{a+b-1}K_{c-2}+K_{a-1}K_{b-1}K_{c-3}>0.
\]
The first difference is nonnegative by \eqref{eq:bare-injection}; a remaining
term is positive.  Therefore $R>Z_q(B-A)$ also in the $0\pmod4$ channel, and
the preceding sign split proves \eqref{eq:one-phase}.
\end{proof}

\begin{lemma}[Two-hostile channel]\label{lem:two}
If both odd cycles are hostile, then
\begin{equation}\label{eq:two-phase}
 \Theta_G(t)\leq
 \Arg(Z_{|C_1|}(t)+2i)+\Arg(Z_{|C_2|}(t)+2i).
\end{equation}
The sum on the right is its continuous product argument, not a principal
argument reset when the real part of the product changes sign.
\end{lemma}

\begin{proof}
Write $z_j=Z_{C_j}(a|_{C_j})$, $B_1=d_2$, and $B_2=d_1$.  The packet identity
with the other path omitted gives
\begin{equation}\label{eq:Ej}
 R=z_jB_j+E_j,
 \qquad E_j=z_j^xL_Q+z_j^yR_Q+z_j^{xy}M_Q+2\varepsilon d_0.
\end{equation}
For $\varepsilon=+1$, clearly $E_j\geq0$.  For $\varepsilon=-1$, the residues
force $P_0$ odd and $P_1,P_2$ even.  Each omitted path $Q$ has even length at
least two, so $L_Q,R_Q\geq1$.  Retaining a matching of
$P_0-\{x,y\}$ and adding the unique perfect matching of the remaining even
path injects coefficient-preservingly into matchings of $C_j-x$ and $C_j-y$.
Consequently the delicate expression factors exactly as
\begin{align}\label{eq:zero-factor}
 E_j={}&(z_j^x-d_0)L_Q+d_0(L_Q-1)\notag\\
       &+(z_j^y-d_0)R_Q+d_0(R_Q-1)+z_j^{xy}M_Q\geq0.
\end{align}
This remains valid when $P_0$ is one edge, when $Q$ has length two, and when
$M_Q=0$; no empty-state convention has been suppressed.

Thus $2B_j/R\leq2/z_j$.  Put $\theta_j=\arctan(2/z_j)$ and
$\Theta=\arctan(2(B_1+B_2)/R)\in(0,\pi/2)$.  If
$\theta_1+\theta_2<\pi/2$, then
\[
 \tan(\theta_1+\theta_2)
 =\frac{2/z_1+2/z_2}{1-4/(z_1z_2)}
 \geq\frac{2(B_1+B_2)}R,
\]
so $\Theta\leq\theta_1+\theta_2$.  If the sum is at least $\pi/2$, the result
is immediate.  Moreover
\[
 (z_1+2i)(z_2+2i)=(z_1z_2-4)+2i(z_1+z_2)
\]
has continuous argument $\theta_1+\theta_2\in(0,\pi)$, passing through
$\pi/2$ rather than jumping branches.  Apply \eqref{eq:weighted-cycle} to
 both cycles to obtain \eqref{eq:two-phase}.
\end{proof}

\section{The phase charge theorem}

\begin{theorem}[Attached-theta phase charge]\label{thm:theta}
Let $G$ be a finite simple connected graph whose $2$-core is a theta, with
arbitrary finite rooted trees attached at arbitrary core vertices.  If $G$ is
nonbipartite, then
\begin{equation}\label{eq:theta-bound}
 D(G)\geq-4(\sqrt5-2),
 \qquad \splus(G)-|V(G)|\geq5-2\sqrt5>0.
\end{equation}
If $G$ is bipartite, then $\splus(G)=|V(G)|+1$.
\end{theorem}

\begin{proof}
For a hostile odd cycle $C_q$, $q\equiv1\pmod4$, its spectrum and
\eqref{eq:coulson} give
\begin{equation}\label{eq:cycle-charge}
 D(C_q)=-2(\sec(\pi/q)-1),
 \qquad
 \int_0^\infty t\Arg(Z_q(t)+2i)\,dt
 =\frac\pi2(\sec(\pi/q)-1).
\end{equation}
For $q\geq5$, this charge decreases and is at most
$\delta:=\sec(\pi/5)-1=\sqrt5-2$.  There are at most two hostile cycles.
The zero-hostile phase is negative, Lemma~\ref{lem:one} charges one hostile
cycle to its bare cycle, and Lemma~\ref{lem:two} charges two to their sum.
Thus $D(G)\geq-4\delta$.  By \eqref{eq:energy},
\[
 \splus(G)-n=1+\frac12D(G)
 \geq1-2\delta=5-2\sqrt5>0.
\]
The last inequality follows from $25>20$.  The bipartite assertion was proved
before \eqref{eq:cycles}.
\end{proof}

\begin{proof}[Proof of Theorem~\ref{thm:main}]
By Proposition~\ref{prop:structure}, a cactus core makes $G$ a rank-two
cactus, covered by Theorem~\ref{thm:cactus}; a theta core is covered by
Theorem~\ref{thm:theta}.  Both alternatives include arbitrary tree attachments.
\end{proof}

\section{Reproducibility and scope}

The full derivation underlying Sections 2--5 is the research note
\cite{ThetaNote}; the proof above reproduces its uniform argument and key
lemmas.  The exact verifier \cite{ThetaVerify} checks the endpoint states,
low-length conventions, bare injection, $0\pmod4$ factorization, phase ledger,
and rejection of seven hostile mutations.  From the repository root, run it
normally and with optimization:
\begin{verbatim}
python3 positive-square-energy/experiments/\
  arbitrary_attached_theta_phase_verifier.py
python3 -O positive-square-energy/experiments/\
  arbitrary_attached_theta_phase_verifier.py
bash scripts/build-paper.sh all-bicyclic-graphs
\end{verbatim}
The verifier is corroborative only: uniform completeness comes from the
continuant recurrences and matching injections, not a finite path census.

\paragraph{Scoped nonclaim.}
The theorem concerns finite simple connected graphs with exactly $n+1$ edges.
It does not claim a strict bound for unicyclic or disconnected graphs, or an
all-densities strengthening.  Indeed $\splus(C_4)=|V(C_4)|=4$.

\section*{AI disclosure}
This manuscript and its proof-object assembly were generated by an AI coding
and mathematical reasoning system from the cited repository materials.  No
human authorship, review, or verification is claimed.  The result should be
assessed from the displayed argument and reproducible dependencies.  No
external publication was performed as part of this work.

\begin{thebibliography}{99}
\bibitem{AKMPZ} S.~Akbari, H.~Kumar, B.~Mohar, S.~Pragada, and S.~Zhang,
\emph{Refinement of a conjecture on positive square energy of graphs},
arXiv:2506.07264, 2025.

\bibitem{Cactus} Companion manuscript, \emph{Positive Square Energy of Every
Nonunicyclic Cactus}, 2026; \path{all-nonunicyclic-cacti/paper.tex}.

\bibitem{ThetaNote} \emph{Arbitrary attached theta: continuant phase
reduction}, research proof note, 2026;
\path{positive-square-energy/bicyclic-theta/arbitrary-attached-theta-phase-note.md}.

\bibitem{ThetaVerify} \emph{Supporting fail-closed checks for the attached-
theta phase proof}, 2026;
\path{positive-square-energy/experiments/arbitrary_attached_theta_phase_verifier.py}.
\end{thebibliography}

\end{document}
